\documentclass[9pt,a4paper]{article}

\usepackage[utf8]{inputenc}
\usepackage[T1]{fontenc}
\usepackage[
  a4paper,
  top=1cm,
  left=1cm,
  right=1cm,
  bottom=1.4cm,
  footskip=0.6cm
]{geometry}
\usepackage{multicol}
\usepackage{booktabs}
\usepackage{longtable}
\usepackage{listings}
\usepackage{enumitem}
\usepackage{titlesec}

\setlength{\parindent}{0pt}
\setlength{\parskip}{2pt}

\titlespacing*{\section}{0pt}{4pt}{2pt}
\titlespacing*{\subsection}{0pt}{3pt}{1pt}
\titlespacing*{\subsubsection}{0pt}{2pt}{1pt}

\lstset{
    basicstyle=\ttfamily\scriptsize,
    frame=single,
    breaklines=true
}

\title{\vspace{-1cm}Snowflake Cloud Data Warehouse Cheat Sheet}
\author{Markaj Agron \and Jorand Yuuta \and Stampfli Nathan \and Liao Pei-Wen}
\date{June 2026}

\begin{document}

\maketitle
\vspace{-1cm}

\section{Why Snowflake?}

\subsection{Real-world Scenario: SwissBike SA}

SwissBike SA is a Swiss SME selling bicycles and accessories online. Data comes from multiple sources: e-commerce platform, ERP system, CRM system and marketing tools.

Management needs answers to questions such as: Which countries generate the most revenue? Who are the most valuable customers? Which products should be promoted? How are sales evolving over time?

Without a centralized solution, data remains scattered across systems, making reporting slow, inconsistent and difficult to maintain.

\subsection{What does Snowflake do?}

Snowflake is a cloud-native data warehouse that centralizes data from multiple sources and enables large-scale analytical processing.

\subsubsection{Data Storage Characteristics}

Supports structured, semi-structured and unstructured data; native support for JSON, Avro, ORC, Parquet and XML; automatic compression; micro-partitioning; separation of Storage and Compute.

\subsubsection{Inputs and Outputs}

\begin{minipage}[t]{0.47\columnwidth}
\textbf{Inputs}\\
CSV files, databases, ERP systems, CRM systems, APIs, streaming data.
\end{minipage}
\hfill
\begin{minipage}[t]{0.47\columnwidth}
\textbf{Outputs}\\
SQL query results, dashboards (Power BI, Tableau, Looker), reports, machine learning and analytics data.
\end{minipage}

\section{Benefits and Limitations}

\subsection{Benefits}

Separation of Storage and Compute; broad data support; high scalability; excellent concurrency; massively parallel processing (MPP); minimal administration; strong security; Time Travel \& Fail-safe; fast analytical performance; secure data sharing; multi-cloud support; automatic cost optimization.

\subsection{Limitations}

Cost; unpredictable spending; vendor lock-in; cloud-only platform; limited infrastructure control; Time Travel storage costs; learning curve; not optimized for OLTP workloads; dependency on cloud providers.

\section{Cost Structure}

\subsection{Storage}

Measured in TB stored per month.

\begin{lstlisting}
1 TB stored data
Cost: approximately $23/month
\end{lstlisting}

\subsection{Compute}

Measured in credits consumed by Virtual Warehouses. Cost depends on warehouse size, execution time and number of clusters.

\section{Monthly Cost Example}

\subsection{SwissBike SA}

Assumptions: 1 TB storage, 10 business users, 10 hours warehouse usage/week, X-Small warehouse, Snowflake AI features enabled.

\begin{center}
\begin{tabular}{lr}
\toprule
Component & Cost \\
\midrule
Compute & \$117.00 \\
CoWork & \$123.08 \\
CoCo & \$5.78 \\
Storage & \$23.00 \\
\midrule
Total & \$268.86 \\
\bottomrule
\end{tabular}
\end{center}

\textbf{Observations:} Compute is the largest cost component. AI features significantly increase monthly spending. Storage costs remain relatively low. Automatic warehouse suspension helps reduce expenses.

\section{Getting Started}

\subsection{Prerequisites}

$\bullet$ Snowflake account \quad
$\bullet$ Access to AWS, Azure or GCP \quad
$\bullet$ Warehouse creation permissions

\subsection{Setup Steps}

1. Create database $\rightarrow$
2. Create schema $\rightarrow$
3. Create warehouse $\rightarrow$
4. Load data $\rightarrow$
5. Execute queries

\subsection{Basic SQL Operations}

Snowflake relies on standard SQL for data definition (DDL), data manipulation (DML) and administration tasks.

\begin{lstlisting}[language=SQL]
CREATE DATABASE POC_SNOWFLAKE;
CREATE SCHEMA DEMO;

CREATE WAREHOUSE DEMO_WH
WITH WAREHOUSE_SIZE='XSMALL';

SELECT CURRENT_ACCOUNT(),
       CURRENT_USER(),
       CURRENT_WAREHOUSE();

SELECT *
FROM SNOWFLAKE_SAMPLE_DATA
.TPCH_SF1.CUSTOMER
LIMIT 10;
\end{lstlisting}
\section{Common Commands}

\textbf{Query sample data}

\begin{lstlisting}[language=SQL]
SELECT *
FROM SNOWFLAKE_SAMPLE_DATA
.TPCH_SF1.ORDERS
LIMIT 100;
\end{lstlisting}

\textbf{Warehouse operations}

\begin{lstlisting}[language=SQL]
ALTER WAREHOUSE COMPUTE_WH
SET WAREHOUSE_SIZE='SMALL';

ALTER WAREHOUSE COMPUTE_WH
SUSPEND;

ALTER WAREHOUSE COMPUTE_WH
RESUME;
\end{lstlisting}

\textbf{Time Travel}

\begin{lstlisting}[language=SQL]
DROP TABLE CUSTOMER_PROD;
UNDROP TABLE CUSTOMER_PROD;
\end{lstlisting}

Time Travel recovers historical versions of data; Fail-safe provides an additional recovery period.

\textbf{Zero-Copy Clone}

\begin{lstlisting}[language=SQL]
CREATE TABLE CUSTOMER_TEST
CLONE CUSTOMER_PROD;
\end{lstlisting}

\textbf{Query Profile}

Monitoring $\rightarrow$ Query History $\rightarrow$ Query Profile.

Displays Table Scan, Join, Aggregate and Sort operations, helping understand MPP execution.

\section{When Should a SME Use Snowflake?}

\textbf{Recommended:} Business Intelligence, Data Warehousing, Analytics, Reporting, Data Science, Multi-source Data Integration.

\textbf{Not Recommended:} Small applications with limited data, OLTP workloads, projects requiring full infrastructure control.

\textbf{Recommendation:} Snowflake is an excellent choice for SMEs requiring scalable analytics with minimal operational overhead, provided costs and vendor lock-in risks are actively monitored.

\end{document}